\DocumentMetadata{
  lang=en-US,
  pdfstandard=ua-2,
  testphase={phase-III,title,table}
}
\documentclass[letterpaper,11pt]{article}

\newcommand{\HomeworkTitle}{Component Showcase}
\newcommand{\HomeworkSubtitle}{Living documentation for homework-template}
\newcommand{\StudentName}{Student Name}
\newcommand{\StudentId}{Student ID}
\newcommand{\CourseCode}{COURSE 101}
\newcommand{\CourseName}{Course Name}
\newcommand{\InstructorName}{Instructor Name}
\newcommand{\SubmissionDate}{Not applicable}
\newcommand{\Collaborators}{None}

\usepackage{theme/ricardo-homework}

\hypersetup{
  pdftitle={Homework Template Showcase},
  pdfauthor={homework-template},
  pdfsubject={Component catalogue for homework-template},
  pdfkeywords={LaTeX, homework, template, showcase}
}

\begin{document}

\makehomeworktitle

\section{Cheat sheet}

Everything this template offers in a quick reference.
Each entry names the environment or command that produces it; the sections that follow show them in use.

\begin{multicols}{2}
\noindent
\textbf{\color{rzaccent}Structure}\par
\code{problem} \dotfill\ a numbered problem\par
\code{subproblems} \dotfill\ (a), (b), (c)\par
\code{solution} \dotfill\ worked reasoning\par
\code{answer} \dotfill\ the final result\par

\medskip
\noindent
\textbf{\color{rzaccent}Callouts}\par
\code{callout} \dotfill\ neutral note\par
\code{tip} \dotfill\ what to remember\par
\code{warn} \dotfill\ a caveat\par
\code{caution} \dotfill\ a mistake to avoid\par

\medskip
\noindent
\textbf{\color{rzaccent}Drafting}\par
\code{\textbackslash HomeworkComment} \dotfill\ review note\par
\code{\textbackslash DeclareCommenter} \dotfill\ author shorthand\par
\code{\textbackslash TODO} \dotfill\ draft marker\par
\code{\textbackslash HIDDEN} \dotfill\ source-only note\par

\medskip
\noindent
\textbf{\color{rzaccent}Mathematics}\par
\code{theorem} \code{lemma} \code{proposition}\par
\code{corollary} \code{definition} \code{remark}\par
\code{proof} \dotfill\ ends with a tombstone\par
\code{smalldescription} \dotfill\ compact definitions\par

\columnbreak
\noindent
\textbf{\color{rzaccent}Code}\par
\code{codebox[lang]\{file\}} \dotfill\ titled plate\par
\code{codeplain[lang]} \dotfill\ no title bar\par
\code{\textbackslash inputcode} \dotfill\ a real file\par
\code{terminal} \dotfill\ shell transcript\par
\code{\textbackslash code} \code{\textbackslash filepath} \dotfill\ inline\par

\medskip
\noindent
\textbf{\color{rzaccent}Figures}\par
\code{\textbackslash HomeworkFigure} \dotfill\ an image\par
\code{\textbackslash HomeworkDiagram} \dotfill\ a D2 diagram\par

\medskip
\noindent
\textbf{\color{rzaccent}Marks}\par
\rzcheck\ \code{\textbackslash rzcheck} \quad
\rzcross\ \code{\textbackslash rzcross}\par
\rzwarnglyph\ \code{\textbackslash rzwarnglyph} \quad
\rzinfoglyph\ \code{\textbackslash rzinfoglyph}\par
\end{multicols}

\begin{callout}[Two files]
  Write the assignment, including its metadata and every problem, in \filepath{homework.tex}.
  Keep \filepath{showcase.tex} as the independent component reference.
\end{callout}

\section{Problems and solutions}

\begin{problem}[12 points]{A numbered problem}
  Problem statements may contain prose, displayed mathematics, lists, and references.
  The optional argument records the available points.
  \begin{subproblems}
    \item The first subproblem.
    \item The second subproblem.
  \end{subproblems}

  \begin{solution}
    Solutions use a distinct breakable panel and may span multiple pages.
  \end{solution}

  \begin{answer}
    Use an answer box when the final result should be easy to locate.
  \end{answer}
\end{problem}

\begin{problem}{A problem without points}
  The points argument is optional, but the descriptive title is required.
\end{problem}

\subsection*{Lists}

Bullets and numbers both take the accent on the marker and leave the text in ink, so a list reads as structure without the words changing colour.

\begin{itemize}
  \item A first point, at the top level.
  \item A second one, with a nested list under it.
  \begin{itemize}
    \item The second level takes a dash.
    \item Only the marker is coloured, at every depth.
    \begin{itemize}
      \item And the third an asterisk.
    \end{itemize}
  \end{itemize}
  \item Back at the top level.
\end{itemize}

\begin{enumerate}
  \item A numbered list matches the lettered \code{subproblems} above.
  \item Both are set once, in \filepath{theme/ricardo-homework.sty}.
\end{enumerate}

\section{Callouts}

Four boxes for interrupting the argument.
Each carries an icon and a word as well as a colour, so a reader in greyscale still gets which kind it is.
The optional argument replaces the label.

\begin{callout}
  Neutral context.
  The default label comes from the environment name.
\end{callout}

\begin{tip}
  The one sentence you want the grader to remember.
\end{tip}

\begin{warn}[Watch the bound]
  The recurrence only closes for $n\geq 1$; the base case has to be checked separately.
\end{warn}

\begin{caution}[A common mistake]
  Do not assume the loop terminates because it does on your inputs.
\end{caution}

\begin{callout}[Custom label]
  All four take an optional argument.
  Giving one turns the label into the sentence it introduces, rather than a generic banner.
\end{callout}

\section{Draft annotations and writing helpers}

Review comments are visible while \code{\textbackslash TemplateStatus} is not
\code{final}. The generic command takes an optional author label:

\HomeworkComment[Instructor]{Check whether this argument covers the empty input.}

Declare a short command once when a reviewer comments repeatedly:

\DeclareCommenter{\Reviewer}{Reviewer}
\Reviewer{State the invariant before using it in the induction step.}

The \TODO{marker identifies wording that still needs attention}. The text remains
when comments are hidden, but the marker does not.
\HIDDEN{This editorial note never appears in any mode.}

\begin{smalldescription}
  \item[Abbreviations] \ie \eg and \etal insert their following space safely.
  \item[Marks] \cmark\ accepted, \xmark\ rejected, with text alternatives.
  \item[Comparisons] A \better{better result} and a \worse{worse result} use
    the theme's checked state colours.
\end{smalldescription}

\section{Mathematics}

\begin{theorem}[Geometric series]
  For $r\neq 1$ and every integer $n\geq 0$,
  \[
    \sum_{k=0}^{n}r^k=\frac{1-r^{n+1}}{1-r}.
  \]
\end{theorem}

\begin{proof}
  Multiply the sum by $1-r$ and cancel adjacent terms.
\end{proof}

The template includes theorem, lemma, proposition, corollary, definition, and remark environments.
Equations, theorems, figures, and tables can be referenced with \verb|\cref|.

\section{Images and diagrams}

\HomeworkFigure[0.58\linewidth]
  {example-image-a}
  {A sample raster-compatible figure included by the TeX distribution.}
  {A pale rectangle marked with the letter A.}
  {fig:sample-image}

Every figure takes a visual description as its fourth argument, separate from the caption.
The caption is a label for a reader who can see the figure; the description is the figure, in words, for one who cannot.
\Cref{fig:sample-image} carries both.

\HomeworkDiagram[0.42\linewidth]
  {workflow}
  {The source-to-submission workflow, drawn in D2.}
  {A downward flow of four boxes: homework.tex goes through lualatex to validation, and validation produces homework.pdf.}
  {fig:workflow}

Diagrams are written in D2 under \filepath{assets/diagrams/} and compiled to PDF by \code{make diagrams}.
The generated PDFs stay committed, so editing a diagram needs D2 installed but compiling the document does not.

A diagram needs its description more than a photograph does, and here for two reasons rather than one.

The first is that the labels are not in the text layer at all.
The pipeline goes through PostScript on its way to PDF, and that hop outlines the glyphs, so \code{pdftotext} returns nothing for a diagram.
That is a deliberate trade: the PostScript route produces a byte-identical PDF on every run, which is what lets the build assert that a committed diagram still matches its source, and the direct route keeps the text but changes bytes between runs.

The second reason would hold even if the words were there.
They would arrive as a loose sequence in drawing order, with no edges, no direction and no shape.
Extractable text is not a description.

\section{Code, tables, and citations}

\begin{codebox}[python]{binary_search.py}
def binary_search(values, target):
    low, high = 0, len(values) - 1
    while low <= high:
        middle = (low + high) // 2
        if values[middle] == target:
            return middle
        if values[middle] < target:
            low = middle + 1
        else:
            high = middle - 1
    return None
\end{codebox}

\par\addvspace{0.8\baselineskip}
\begin{table}[H]
  \centering
  \caption{Asymptotic costs of common search operations.}
  \label{tab:complexity}
  \begingroup
  \tagpdfsetup{table/header-rows={1}}
  \begin{tabular}{@{}lll@{}}
    \textbf{Method} & \textbf{Prerequisite} & \textbf{Worst case} \\[0.35em]
    Linear search & None & $\Theta(n)$ \\
    Binary search & Sorted input & $\Theta(\log n)$ \\
    Hash lookup & Suitable hash table & $\Theta(n)$ \\
  \end{tabular}
  \endgroup
\end{table}
\par\addvspace{0.8\baselineskip}

\Cref{tab:complexity} uses spacing rather than decorative rules or vertical separators.
Bibliographic citations use BibTeX and Natbib; for example, see \citet{knuth1984literate}.

\subsection{Other languages}

The optional argument is the Pygments lexer, so any language Pygments knows is available.
The signature is the slides project's on purpose: content written for one template reads the same in the other.

\begin{codebox}[c]{sum_array.c}
size_t sum_array(const int *values, size_t n) {
    size_t total = 0;
    for (size_t i = 0; i < n; i++) {
        total += values[i];  /* no overflow check, on purpose */
    }
    return total;
}
\end{codebox}

\begin{codebox}[haskell]{Fib.hs}
fib :: Int -> Integer
fib n = fibs !! n
  where fibs = 0 : 1 : zipWith (+) fibs (tail fibs)
\end{codebox}

\subsection{Without a title bar}

\code{codeplain} drops the tab and the line numbers, for an excerpt short enough that neither earns its space.

\begin{codeplain}[sql]
SELECT course, AVG(grade) AS mean
  FROM submissions
 GROUP BY course
HAVING COUNT(*) > 5;
\end{codeplain}

\subsection{Inline}

Set an identifier with \code{\textbackslash code}, as in \code{binary\_search(values, target)}, and a path with \code{\textbackslash filepath}, as in \filepath{homework.tex}.
Both stay in Ubuntu Mono, so an identifier reads the same inline as it does inside a plate.

\section{Terminal transcripts}

When the answer is what a command printed, quote the session rather than paraphrasing it.
The body is verbatim, so nothing inside needs escaping except the arguments of the macros themselves.

\TermHost{arch}
\TermPath{~/homework}
\begin{terminal}
\cmd{make check}
\dimout{>> homework: LaTeX pass 1}
\out{homework.pdf: US Letter, tagged}
\ok{all 42 pairs meet WCAG 2.1 AA}
\cmdfail{make submit}
\bad{permission denied}
\warnout{2 problems still marked TODO}
\end{terminal}

The prompt follows the \code{ginger} zsh theme: a green arrow, then the host, then the directory.
\code{\textbackslash cmdfail} turns the arrow red \emph{and} changes it to a cross, so a failed command is not distinguished by colour alone.
\code{\textbackslash TermHost} and \code{\textbackslash TermPath} set the two labels, and both read their argument verbatim.
That matters: a shell path is made of the characters LaTeX reserves, so \code{\textbackslash TermPath\{\textasciitilde/src/my\_app\}} is typed the way the shell shows it rather than escaped.
The path defaults to \code{\textasciitilde}; pass an empty argument to drop it.

\bibliographystyle{plainnat}
\bibliography{references}

\end{document}
